\documentclass[10pt]{extarticle}
\usepackage{apMod}
\usepackage{hyperref}
% Licensed: https://www.gnu.org/licenses/gpl-3.0.html
% Template author: Anindya Prithvi
\vspace{-20pt}
% === Personal Information ===
\newcommand{\yourname}{Himanshu Kumar}%
\newcommand{\youremail}{himanshu22215@iiitd.ac.in}% <-- Update to your real email
\newcommand{\youraddress}{91-8368612369 \\ \href{https://www.linkedin.com/in/himanshu-kumar-95626634b/}{LinkedIn} }%\quad \href{https://github.com/Himanshu8368}{GitHub}}% <-- Update links
\newcommand{\yourrollno}{2022215}
\newcommand{\yourDOB}{09-03-2004}% Optional
\newcommand{\yourgithub}{https://github.com/Himanshu8368}% <-- Update
\newcommand{\yourphoto}{myphoto}% Only the filename without extension
\newcommand{\contain}{0}% Leave as is





\begin{document}
\cvheader

% === Education Section ===
\begin{contained}
\sectionC{Education}
\begin{longtable}{p{0.8\textwidth}p{0.19\textwidth}p{0.01\textwidth}}
    \cveducationitem{IIIT Delhi}{B.Tech (Computer Science and Engineering)}{2022 - 2026 (Present)}{CGPA: 6.33 (Till 6\textsuperscript{th} semester)}
    \cveducationitem{RSBV Patparganj, Delhi-91}{CBSE, Standard 12}{2019 - 2020}{Percentage: 84\%}
    \cveducationitem{RSBV School Block, Shakarpur, Delhi-91}{CBSE, Standard 10}{2017 - 2018}{Percentage: 82\%}
\end{longtable}
\end{contained}

% === Skills Section ===
\begin{contained}
\sectionC{Skills}
\begin{longtable}{p{0.2\textwidth}p{0.75\textwidth}p{0.05\textwidth}}
    \textbf{Expertise Area:}&
    Data Structures and Algorithms (DSA)\\
    \textbf{Programming Languages:}&
    Python, Java, C, C++, Kotlin\\
    \textbf{Tools and Technologies:}&
    Linux, JavaFX, Android Studio, VS Code, IntelliJ IDEA\\
    \textbf{Technical Electives:}&
    Operating System, Database Management Systems (DBMS), Algorithm Design Analysis (ADA), Mobile Computing\\
\end{longtable}
\end{contained}

% === Projects Section ===
\begin{contained}
\sectionC{Projects}
\begin{longtable}{p{0.8\textwidth}p{0.19\textwidth}p{0.01\textwidth}}
    
    \textbf{Plant Diseases Detection with Solution} \newline
    Guide: Arani Bhattacharya \newline
    Developed a mobile application for plant disease detection using ResNet-50 with 96\% accuracy on 38 classes. Used PyTorch Mobile for offline inference. Built advisory system using local LLMs (Llama-2/phi-2). Created Android app with Jetpack Compose, CameraX, and Navigation. 
    \newline \textbf{Tech Stack:} Kotlin, Android Studio, PyTorch Mobile, OpenAI SDK. \newline
    \href{https://github.com/Himanshu8368/MC_project}{GitHub} &
    Team Size: 1 \newline (Jan 2024 - May 2025)\\


    \textbf{Gendered Abuse Detection in Indic Languages} \newline
    Collaborators: Aman Tomar, Deepesh \newline
    Developed machine learning models for detecting gender-based abuse in Hindi, Tamil, and Indian English social media content. Implemented CNN-BiLSTM-Attention architecture with XLM-RoBERTa embeddings, achieving superior performance over IndicBERT baselines. Built multi-task learning system for joint prediction of gendered abuse and explicit content detection. Applied transfer learning using domain-adapted Twitter models for cross-lingual abuse detection.
    \newline \textbf{Tech Stack:} Python, PyTorch, XLM-RoBERTa, CNN, BiLSTM, Attention Mechanisms. \newline
    Team Size: 3 \\

    \textbf{AgriPredict: Machine Learning-Based Crop Yield Prediction} \newline
    Collaborators: Naman Jindal, Nishchay Sharma, Harshil Handoo \newline
    Developed comprehensive ML system for crop yield prediction using multi-country agricultural dataset (28,242 records). Implemented ensemble methods including Random Forest, XGBoost, and Bagging Regressor achieving 98.9\% accuracy. Created interactive GUI using Streamlit for model selection and visualization. Applied advanced data preprocessing, feature engineering with one-hot encoding, and k-fold cross-validation for robust model evaluation.
    \newline \textbf{Tech Stack:} Python, Scikit-learn, XGBoost, Streamlit, Pandas, Matplotlib. \newline
    Team Size: 4 \\

    \textbf{Stick Hero Game} \newline
    Guide: Raghava Mutharaju \newline
    Developed interactive Stick Hero game using JavaFX. Implemented game mechanics and GUI in Java.
    \newline \href{https://github.com/Himanshu8368/Ap-Project}{GitHub} &
    Team Size: 2 \newline (Nov 2023 – Dec 2023)\\

    \textbf{Zoo Management System} \newline
    Guide: Raghava Mutharaju \newline
    Designed a zoo management system in Java as part of Advanced Programming course.
    \newline \href{https://github.com/yourusername/zoo-management}{GitHub} &
    Team Size: 1 \newline (Oct 2023 – Nov 2023)\\

    % \textbf{Simple and Smart Loader} \newline
    % Guide: Vivek Kumar \newline
    % Implemented system-level loaders in C. Developed efficient mechanisms for basic and smart loading operations.
    % \newline \href{https://github.com/yourusername/loader}{GitHub} &
    % Team Size: 2 \newline (Apr 2023 – Jun 2023)\\
\end{longtable}
\end{contained}

% === Interests Section ===
\begin{contained}
\sectionC{Interests and Hobbies}
\begin{itemize}
    \item Playing Table Tennis and Cricket
    \item Geopolitics
\end{itemize}
\end{contained}
\vspace{10pt}
% === Declaration ===
Declaration: The above information is correct to the best of my knowledge.\newline
Himanshu Kumar\newline
Date: 08/07/2025

\end{document}
